% applications/vehicle_space_and_platforms-DE.tex

\chapter{Fahrzeugraum und Bahnsteiginteraktion}
\label{chap:applications-vehicle-space-platform}

Dieses Kapitel wendet die kanonische pose3- und Laufzeitinterpretation auf
Anwendungsfälle des Fahrzeugraums und der Bahnsteigkante an.

\section{Anwendungskontext}

Bahnsteiginteraktion und Lichtraumbewertung hängen vom Verhalten des
Fahrzeugraums relativ zum Laufzeitzustand der Eisenbahn ab. Das relevante
Ingenieurobjekt ist daher nicht allein die Mittelliniengeometrie, sondern
der aus kanonischen Laufzeitdiensten ausgewertete, abgeleitete
fahrzeugrelative Zustand.

\section{Abgeleitete Fahrzeugrauminterpretation}

Das Fahrzeugraumverhalten wird der kanonischen Auswertung von pose3,
Frame und Überhöhung nachgelagert interpretiert. Es ist eine
Anwendungsinterpretation und keine neue Ebene des Zustandsmodells.

Dies ermöglicht die kontextspezifische Bewertung von Bahnsteigspalten,
Lichtraum und Hüllkurvenverhalten unter Betriebsbedingungen.

\section{Ingenieurwissenschaftliche Anwendungsfälle}

Typische Anwendungsfälle sind die Prüfung von Bahnsteigkanten,
Hüllkurvenprüfungen in Bögen und szenarienbasierte Bewertungen für
unterschiedliche Fahrzeugklassen. Sie verwenden die kanonischen Laufzeit-
und Realisierungsebenen.

\section{Repräsentationsgrenze}

Gerenderte oder exportierte Geometrie gilt als nachgelagerte
Repräsentation eines zur Laufzeit ausgewerteten Zustands und nicht für
sich allein als ingenieurwissenschaftliche Wahrheit.

\section{Verbindung zu AIM}

\begin{summary}
Fahrzeugraum- und Bahnsteiganwendungen in AIM sind nachgelagerte
Interpretationen, die auf kanonischen pose3- und Laufzeitdiensten beruhen.

Sie zeigen praktische Entscheidungsunterstützung und wahren zugleich die
Trennung zwischen Ingenieurobjekten und Repräsentationen.
\end{summary}
